\chapter{Einleitung}

Diese Vorlage dient als Grundlage für die formale und inhaltliche Gestaltung wissenschaftlicher Abschlussarbeiten an der Professur für Wirtschaftsinformatik der Freien Universität Berlin. Die Einleitung einer wissenschaftlichen Arbeit motiviert das Thema, grenzt den Untersuchungsgegenstand ab und macht die Relevanz der Forschungsfrage deutlich. Sie bereitet die Lesenden auf die nachfolgenden Kapitel vor und schafft ein klares Verständnis für Zielsetzung und Aufbau der Arbeit.
Hinweis zur Formatierung: Seitenzahlen, Seitenränder, Schrift und Überschriften sind in der \texttt{main.tex} bereits eingestellt und müssen nicht angepasst werden. Achten Sie darauf, dass keine Leerräume bzw. unmotivierten Leerzeilen in der Arbeit vorkommen.
Zwischen einer Kapitelüberschrift und der ersten Unterüberschrift sollte immer ein einleitender Text stehen, der in das Kapitel einführt und dessen Inhalt und Struktur kurz beschreibt. Vermeiden Sie es, direkt von einer Überschrift zur nächsten zu springen, ohne den Lesenden Orientierung zu geben.


\section{Problemumfeld}

Im Abschnitt \glqq{}Problemumfeld\grqq{} wird das wissenschaftliche oder praktische Problem beschrieben, das die Arbeit adressiert. Hier soll deutlich werden, warum das gewählte Thema relevant ist und in welchem Kontext es steht. Eine klare Einordnung in den Forschungsstand zeigt, welche Lücke die Arbeit schließt oder welche neue Perspektive sie eröffnet.
Die Darstellung sollte nachvollziehbar machen, worin die Herausforderung besteht und warum diese bearbeitet werden muss. Eine präzise Beschreibung des Problemumfelds bildet die Grundlage für eine klare Zielsetzung und Forschungsfrage.

In ac felis quis tortor malesuada pretium. Pellentesque auctor neque nec urna. Proin sapien ipsum, porta a, auctor quis, euismod ut, mi. Aenean viverra rhoncus pede. Pellentesque habitant morbi tristique senectus et netus et malesuada fames ac turpis egestas. Ut non enim eleifend felis pretium feugiat. Vivamus quis mi. Phasellus a est. Phasellus magna. In hac habitasse platea dictumst. Curabitur at lacus ac velit ornare lobortis. Curabitur a felis in nunc fringilla tristique. 

\section{Zielsetzung}

In diesem Abschnitt wird die Zielsetzung der Arbeit formuliert. Hier sollte präzise dargelegt werden, was die Arbeit erreichen soll: Welche Forschungsfrage wird beantwortet? Welche Hypothesen werden untersucht? Welche Erkenntnisse sollen gewonnen werden?
Die Zielsetzung sollte klar von der bloßen Themenbeschreibung abgegrenzt sein. Falls die Arbeit bewusst eingegrenzt wird – etwa auf bestimmte Methoden, Branchen oder theoretische Perspektiven – sollte diese Abgrenzung hier transparent gemacht werden. Eine präzise Zielsetzung erleichtert es den Lesenden, die nachfolgende Argumentation nachzuvollziehen und die Ergebnisse später einzuordnen.


\section{Aufbau der Arbeit}

Dieser Abschnitt gibt einen Überblick über die Struktur der Arbeit und beschreibt die Funktion der nachfolgenden Kapitel. Typischerweise umfasst eine wissenschaftliche Arbeit Kapitel zu Grundlagen, Methodik, Ergebnissen und Fazit.
Ein Musteraufbau könnte wie folgt aussehen: Nach der Einleitung werden in Kapitel 2 die theoretischen Grundlagen und der Stand der Forschung dargestellt. Kapitel 3 erläutert die verwendete Methodik. Kapitel 4 präsentiert die Ergebnisse der Untersuchung. Kapitel 5 fasst die Erkenntnisse zusammen, diskutiert die Grenzen der Arbeit und gibt einen Ausblick auf weiterführende Forschung.
Durch die frühzeitige Beschreibung der Argumentationslinie wird die Struktur der Arbeit transparent und für die Lesenden nachvollziehbar.


